\exercisesfor{6}

\begin{enumerate}
\def\labelenumi{\arabic{enumi}.}
\item
  Let D be a set of elements of \(\mathbf{SL}(2,\mathbf{C})\) of the form \(\begin{pmatrix}a & b \\ 0 & a^{-1}\end{pmatrix}\) where a \textgreater{} 0 and \(b \in C\) . Show that D is a Lie subgroup of the underlying real Lie group of \(\mathbf{SL}(2,\mathbf{C})\) and that the mapping \((u,d) \mapsto ud\) of \(\mathbf{SU}(2,\mathbf{C}) \times \mathbf{D}\) into \(\mathbf{SL}(2,\mathbf{C})\) is an isomorphism of real analytic manifolds. Deduce from this and §3, Exercise 7 (c) that \(\mathbf{SL}(2,\mathbf{C})\) is simply connected.
\item
  Let \(G\) be the universal covering of \(\mathbf{SL}(2, \mathbf{R})\) , \(\pi\) the canonical morphism of \(G\) onto \(\mathbf{SL}(2, \mathbf{R})\) and \(N\) its kernel.
\end{enumerate}

\begin{enumerate}
\def\labelenumi{(\alph{enumi})}
\item
  Arguing as in Exercise 1, show that there exists an isomorphism of the real analytic manifold \(\mathbf{U} \times \mathbf{R}^2\) onto the real analytic manifold \(\mathbf{SL}(2, \mathbf{R})\) . Deduce that \(N\) is isomorphic to \(\mathbf{Z}\) .
\item
  If \(\rho\) is an analytic linear representation of \(G\) on a complex vector space, then \(\operatorname{Ker} \rho \supset N\) . (The representation \(\mathrm{L}(\rho)\) of \(\mathfrak{sl}(2, \mathbf{R})\) defines, by complexification, a representation of \(\mathfrak{sl}(2, \mathbf{C})\) and the latter is, by Exercise 1, of the form \(\mathrm{L}(\sigma)\) , where \(\sigma\) is an analytic linear representation of \(\mathbf{SL}(2, \mathbf{C})\) . Then \(\mathrm{L}(\sigma \circ \pi) = \mathrm{L}(\rho)\) and hence \(\sigma \circ \pi = \rho\) .)
\end{enumerate}

\begin{enumerate}
\def\labelenumi{\arabic{enumi}.}
\setcounter{enumi}{2}
\item
  Let \(G_{1}\) be the simply connected nilpotent real Lie group defined in Exercise 5 of § 4. Let Z be the centre of \(G_{1}\) , N a non-trivial discrete subgroup of Z and \(G = G_{1}/N\) . If \(\rho\) is a finite-dimensional analytic linear representation of G, \(\rho(Z/N)\) is a semi-simple set of automorphisms, for Z/N is compact; but \(Z = (G, G)\) and hence \(\rho(Z/N)\) is unipotent (Chapter I, § 6, Proposition 6). Therefore \(\rho\) is trivial on Z/N.
\item
  Let G be a compact connected complex Lie group. Show that every analytic linear representation of G is trivial.
\item
  In the notation of Exercise 9 of § 1, show that \(\mathbf{H}'\) is an integral subgroup of H. Deduce that in the simply connected group \(\mathbf{SU}(2,\mathbf{C})\times \mathbf{SU}(2,\mathbf{C})\) (§ 3, Exercise 7 (c)), there exist one-parameter subgroups which are not closed.
\end{enumerate}

¶ 6. Let I be the interval \(\{1,2\}\) with the discrete topology. Let E be the complete normed space of functions \(f: I \to R\) such that \(\|f\| = \sum_{i \in I} |f(i)| < +\infty\) . For all \(x \in I\) , let \(\varepsilon_x\) be the element of E such that \(\varepsilon_x(t) = 0\) for \(t \neq x\) and \(\varepsilon_x(x) = 1\) . Let P be the subgroup of E generated by the \(x\varepsilon_x\) with \(x \in I\) ; it is a discrete subgroup of E. Let G be the Lie group E/P. Let F be the hyperplane of E consisting of the \(f \in E\) such that \(\sum_{i \in I} f(i) = 0\) . Let H be the Lie group F/(F ∩ P). Let \(\phi\) be the canonical morphism of H into G. Show that \(\phi\) is bijective and is an immersion but is not a Lie group isomorphism. (Let f be a linear form on E with kernel F; show that \(f(P) = R\) and hence that \(F + P = E\) .) Deduce that in Proposition 3 the countability hypothesis cannot be omitted.

\begin{enumerate}
\def\labelenumi{\arabic{enumi}.}
\setcounter{enumi}{6}
\tightlist
\item
  Let \(\phi\) be the morphism of \(\mathbf{Z} / 2\mathbf{Z}\) into the permutation group of \(\mathbf{C}\) which maps the non-identity element to \(z\mapsto \overline{z}\) . Let \(\mathrm{G}\) be the semi-direct product of \(\mathbf{Z} / 2\mathbf{Z}\) by the real Lie group \(\mathbf{C}\) corresponding to \(\phi\) . It is a real Lie group with lie algebra \(\mathbf{R}^2\) . Show that there exists on \(\mathrm{G}\) no complex Lie group structure compatible with the real Lie group structure.
\end{enumerate}

¶ 8. Let H be a complex Hilbert space of dimension \(\aleph_{0}\) and G the unitary group of H considered as a real Lie group. The group G is simply connected. \(\dagger\) (a) Let Z be the centre of G, which is isomorphic to T. Let a be an irrational number. Let \(\hat{\mathbf{a}}\) be the Lie subalgebra of \(\mathrm{L}(\mathrm{G})\times \mathrm{L}(\mathrm{G})\) consisting of the \((x,ax)\) where \(x\in \mathrm{L}(Z)\) . Let S be the corresponding integral subgroup of \(\mathrm{G}\times \mathrm{G}\) . Then \(\hat{\mathbf{a}}\) is an ideal of \(\mathrm{L}(\mathrm{G})\times \mathrm{L}(\mathrm{G})\) ; however S is not closed and is dense in \(\mathbf{Z}\times \mathbf{Z}\) .

\begin{enumerate}
\def\labelenumi{(\alph{enumi})}
\setcounter{enumi}{1}
\tightlist
\item
  Let \(\mathfrak{g} = (\mathrm{L}(\mathrm{G})\times \mathrm{L}(\mathrm{G})) / \tilde{\mathfrak{s}}\) . There exists no Lie group with Lie algebra \(\mathfrak{g}\) . (Let \(\mathrm{H}\) be such a group. As \(\mathrm{G}\) is simply connected, there exists a morphism \(\phi : \mathrm{G}\times \mathrm{G}\to \mathrm{H}\) such that \(\mathrm{L}(\phi)\) is the canonical mapping of \(\mathrm{L}(\mathrm{G})\times \mathrm{L}(\mathrm{G})\) onto \(\mathfrak{g}\) . Let \(\mathrm{N} = \operatorname{Ker}\phi\) . Then \(\mathrm{L}(\mathrm{N})\supset \tilde{\mathfrak{s}}\) , hence \(\mathrm{N}\supset \mathrm{S}\) and therefore \(\mathrm{N}\supset \mathrm{Z}\times \mathrm{Z}\) by (a). Then \(\mathrm{L}(\mathrm{N})\supset \mathrm{L}(\mathrm{Z})\times \mathrm{L}(\mathrm{Z})\) , which is absurd.)
\end{enumerate}

\begin{enumerate}
\def\labelenumi{\arabic{enumi}.}
\setcounter{enumi}{8}
\item
  Let X be a non-empty connected compact complex manifold of dimension n.~Suppose that there exist holomorphic vector fields \(\xi_{1},\ldots,\xi_{n}\) on X which are linearly independent at each point of X. Show that there exist a complex Lie group G and a discrete subgroup D of G such that X is diffeomorphic to G/D. (Let \(c_{ijk}\) be the holomorphic functions on X such that \([\xi_{i},\xi_{j}]=\sum_{k}c_{ijk}\xi_{k}\) . The \(c_{ijk}\) are constants because X is compact. Take G to be the simply connected complex Lie group whose Lie algebra admits the \(c_{ijk}\) as constants of structure and use Theorem 5.)
\item
  Let G be a Lie group with a finite number of connected components. For G to be unimodular, it is necessary and sufficient that Tr ad \(a = 0\) for all \(a \in \mathrm{L}(\mathrm{G})\) .
\item
  Let E be a complete normable space over \(\mathbf{C}\) , \(v \in \mathcal{L}(\mathrm{E})\) and \(g = \exp(v)\) . Suppose that \(\operatorname{Sp}(v) \cap 2i\pi(\mathbf{Z} - \{0\}) = \emptyset\) . Let \(\mathrm{E}_1, \mathrm{E}_2\) be closed vector subspaces of E which are stable under \(v\) , such that \(\mathrm{E}_2 \subset \mathrm{E}_1\) . Suppose that the automorphism of \(\mathrm{E}_1 / \mathrm{E}_2\) derived from \(g\) is the identity. Then \(v(\mathrm{E}_1) \subset \mathrm{E}_2\) .
\end{enumerate}

¶ 12. Consider a real or complex complete normable space E and a closed vector subspace F of E such that \(F^{0}\) admits no topological supplement in \(E^{\prime}\dagger\) (a) Let A (resp. B) be the complete normable space of continuous endomorphisms of E (resp. F). Let C be the set of \(u \in A\) such that \(u(F) \subset F\) . Let \(\alpha\) be the mapping \(u \mapsto (u, u \mid F)\) of C into \(A \times B\) . Then \(\alpha\) is an isomorphism of the complete normable space C onto a closed vector subspace of \(A \times B\) and \(\alpha(C)\) admits no topological supplement in \(A \times B\) . (Let \(x \in E\) be such that \(x \notin F\) . Let \(\xi \in \mathrm{F}^0\) be such that \(\langle x, \xi \rangle = 1\) . If \(\eta \in \mathrm{E}'\) , let \(\tilde{\eta}\) denote the element \(y \mapsto \langle y, \eta \rangle x\) of A. Suppose that there exists a projector \(\pi\) of A × B onto \(\alpha(\mathrm{C})\) . We define \(\tilde{\pi} \colon \mathrm{E}' \to \mathrm{E}'\) by \(\tilde{\pi}(\eta) = t(\alpha^{-1}(\pi(\tilde{\eta}, 0)))(\xi)\) . Then \(\tilde{\pi}\) is a projector of \(\mathrm{E}'\) onto \(\mathrm{F}^0\) , which is absurd.)

as follows: \(M = \bigoplus_{i=0}^{n} M_{i}\) is graded by the \(M_{i}\) ; \(M_{0}\) is the set of scalar multiples of 1; \(M_{i}\) is the set of scalar multiples of a non-zero element u; \(M_{2} = \{0\}\) ; \(M_{3} = F\) ; \(M_{4} = E\) ; \(M_{i} = \{0\}\) for i \textgreater{} 5; ux = x for all \(x \in M_{3}\) . Let N be the complete normable space of continuous endomorphisms of the complete normable space M. Let \(N_{1}\) be the set of continuous derivations of M. Show, using (a), that \(N_{1}\) has no topological supplement in N.

\begin{enumerate}
\def\labelenumi{(\alph{enumi})}
\setcounter{enumi}{2}
\tightlist
\item
  Show that the group of bicontinuous automorphisms of M is a Lie quasi-subgroup of \(\mathbf{GL}(\mathbf{M})\) , but not a Lie subgroup of \(\mathbf{GL}(\mathbf{M})\) . (Use Corollary 2 to Proposition 18 and (b).)
\end{enumerate}

\begin{enumerate}
\def\labelenumi{\arabic{enumi}.}
\setcounter{enumi}{12}
\item
  Let \(\mathbf{G}\) be a real Lie group, \(\mathbf{L} = \mathbf{L}(\mathbf{G})\) , and \(\phi: \mathbf{L} \to \mathbf{G}\) a mapping which is differentiable at 0, such that \(\mathrm{T}_0(\phi) = \mathrm{Id}_{\mathrm{L}}\) and such that \(\phi(nx) = \phi(x)^n\) for all \(x \in \mathbf{L}\) and \(n \in \mathbf{Z}\) . Then \(\phi = \exp_0\) . (Let \(\mathbf{V}\) be a neighbourhood of 0 in \(\mathbf{L}\) and \(\mathbf{W}\) a neighbourhood of \(e\) in \(\mathbf{G}\) such that \(0 = \exp_0 | \mathbf{V}\) is an analytic isomorphism of \(\mathbf{V}\) onto \(\mathbf{W}\) . Let \(\psi = \theta^{-1} \circ (\phi | \phi^{-1}(\mathbf{W}))\) . Then \(\mathrm{T}_0(\psi) = \mathrm{Id}_{\mathrm{L}}\) , whence \(\psi = \mathrm{Id}_{\mathrm{L}}\) using the equality \(\psi\left(\frac{1}{n}x\right) = \frac{1}{n}\psi(x)\) for \(x\) sufficiently close to 0 and \(n \in \mathbf{Z} - \{0\}\) .
\item
  Let \(\mathfrak{a}_1\) be the commutative Lie algebra \(\mathbf{R}^4\) , identified with \(\mathbf{C}^2\) . Let \(\mathfrak{a}_2\) be the commutative Lie algebra \(\mathbf{R}\) . Let \(\phi\) be the homomorphism of \(\mathfrak{a}_2\) into \(\mathrm{Der}(\mathfrak{a}_1)\) which maps 1 to the derivation \((z_1, z_2) \mapsto (iz_1, i\sqrt{2}z_2)\) . Let \(\mathfrak{a}\) be the semi-direct product of \(\mathfrak{a}_2\) by \(\mathfrak{a}_1\) corresponding to \(\phi\) .
\end{enumerate}

\begin{enumerate}
\def\labelenumi{(\alph{enumi})}
\tightlist
\item
  Show that (Int a) \textbar{} \(a_{1}\) contains, for all \(\phi \in R\) , the automorphism
\end{enumerate}

\[
\left(z _ {1}, z _ {2}\right) \mapsto \left(e ^ {i \phi} z _ {1}, e ^ {i \sqrt {2} \phi} z _ {2}\right).
\]

\begin{enumerate}
\def\labelenumi{(\alph{enumi})}
\setcounter{enumi}{1}
\item
  Show that the closure G of (Int a) \textbar{} a₁ in GL(a₁) contains, for all (φ, φ') ∈ R², the automorphism (z₁, z₂) ↦ (eⁱΦz₁, eⁱΦ'z₂).
\item
  Show that \(\mathbf{L}(\mathbf{G})\) contains the endomorphism \(u: (z_1, z_2) \mapsto (z_1, 0)\) of \(\mathfrak{a}_1\) and that there exists no \(x \in \mathfrak{a}\) such that \(u = (\operatorname{ad} x) \mid \mathfrak{a}_1\) .
\item
  Deduce that \(\operatorname{Int}(\mathfrak{a})\) is not closed in \(\operatorname{Aut}(\mathfrak{a})\) .
\end{enumerate}

\begin{enumerate}
\def\labelenumi{\arabic{enumi}.}
\setcounter{enumi}{14}
\tightlist
\item
  Show that the finite-dimensional simply connected real Lie group of §3, Exercise 7 (a) contains non-simply connected Lie subgroups (it contains in fact Lie groups isomorphic to U).
\end{enumerate}

  ¶ 16. (a) Let g be a complex Lie algebra. Let \(\bar{g}\) be the complex Lie algebra derived from g using the automorphism \(\lambda \mapsto \bar{\lambda}\) of C. Let \(g_0\) be the real Lie algebra
derived from g by restriction of scalars. Let \(g' \supset g_0\) be the complex Lie algebra \(g_0 \otimes_{R} C\) . For all \(x \in g\) , we write

\[
f (x) = \frac {1}{2} (x \otimes 1 - (i x) \otimes i) \in \mathfrak {g} ^ {\prime}, \quad g (x) = \frac {1}{2} (x \otimes 1 + (i x) \otimes i) \in \mathfrak {g} ^ {\prime}.
\]

Show that \(f(\text{resp. } g)\) is an isomorphism of \(\mathfrak{g}\) (resp. \(\overline{\mathfrak{g}}\) ) onto an ideal \(\mathfrak{m}\) (resp. \(\mathfrak{n}\) ) of \(\mathfrak{g}'\) and that the ideals \(\mathfrak{m}, \mathfrak{n}\) are supplementary in \(\mathfrak{g}'\) . This defines projectors of \(\mathfrak{g}'\) onto \(\mathfrak{m}, \mathfrak{n}\) . Show that, for all \(x \in \mathfrak{g}, f(x) = p(x), g(x) = q(x)\) .

\begin{enumerate}
\def\labelenumi{(\alph{enumi})}
\setcounter{enumi}{1}
\item
  Suppose that g is finite-dimensional. Let G be the simply connected complex Lie group with Lie algebra g. Let \(\overline{G}\) be the conjugate complex Lie group and \(G_{0}\) the underlying real Lie group. Then \(\mathrm{L}(\overline{\mathbf{G}})=\overline{\mathfrak{g}}\) and \(\mathrm{L}(G_{0})=g_{0}\) . Let M (resp. N) be the simply connected complex Lie group with Lie algebra m (resp. n). Then f defines an isomorphism \(\phi\) of G onto M, g defines an isomorphism \(\gamma\) of \(\overline{G}\) onto N and the simply connected complex Lie group \(G'\) with with Lie algebra \(g'\) is identified with M × N. Show that the real integral subgroup of \(G'\) with Lie algebra \(g_{0}\) is closed and simply connected and is identified with \(G_{0}\) . The restriction to G ⊆ M × N of \(pr_{1}\) (resp. \(pr_{2}\) ) is the isomorphism \(\phi\) (resp. \(\gamma\) ) of G (resp. \(\overline{G}\) ) onto M (resp. N).
\item
  Deduce from (b) and Remark 2 of no. 10 that \(\mathbf{G}'\) , with the canonical injection of \(\mathbf{G}_0\) into \(\mathbf{G}'\) , is the universal complexification of \(\mathbf{G}_0\) .
\item
  Let H be a connected complex Lie group with Lie algebra g, \(\bar{H}\) the conjugate complex Lie group and \(H_{0}\) the underlying real Lie group, such that G is the universal covering space of H. Let K be the (discrete) kernel of the canonical morphism of G onto H. Show that K is a central subgroup of \(G'\) . The canonical injection of \(G_{0}\) into \(G'\) therefore defines an injection i of \(H_{0}\) into \(G'/N\) . Show that \((G'/N, i)\) is the universal complexification of \(H_{0}\) . (Note that this ordered pair has the universal property of Proposition 20.)
\item
  We write \(G^{\prime}/N = (H_{0})_{\mathbf{c}}\) . Deduce from (d) a canonical morphism \(\psi\) of \((\mathrm{H}_{0})_{\mathbf{c}}\) onto \(H \times \bar{H}\) , which is a covering mapping. Show, by the example \(H = C^{*}\) , that \(\psi\) is not in general an isomorphism.
\end{enumerate}

\begin{enumerate}
\def\labelenumi{\arabic{enumi}.}
\setcounter{enumi}{16}
\tightlist
\item
  \begin{enumerate}
  \def\labelenumii{(\alph{enumii})}
  \tightlist
  \item
    Let G, \(\overline{G}\) , \(G_{0}\) be as in Exercise 16 (b). Let V be a finite-dimensional complex vector space and \(\rho\) an irreducible linear representation of \(G_{0}\) on V. Show that there exist finite-dimensional complex vector spaces X, Y, an analytic irreducible linear representation \(\sigma\) (resp. \(\tau\) ) of G (resp. \(\overline{G}\) ) on X (resp. Y) and an isomorphism of V onto X \(\otimes\) Y which transforms \(\rho\) into \(\sigma \otimes \tau\) . (Use Exercise 16 (a), Proposition 2 of Chapter I, § 2, and the Corollary to Proposition 8 of Algebra, Chapter VIII, § 7.)
  \end{enumerate}
\end{enumerate}

\begin{enumerate}
\def\labelenumi{(\alph{enumi})}
\setcounter{enumi}{1}
\tightlist
\item
  Show that the conclusion of (a) is not necessarily true if we omit the hypothesis that G is simply connected. (Consider the complex Lie group \(\mathbf{C}^*\) .)
\end{enumerate}

\begin{enumerate}
\def\labelenumi{\arabic{enumi}.}
\setcounter{enumi}{17}
\tightlist
\item
  Let \(\Lambda\) be a finite-dimensional connected commutative complex Lie group with Lie algebra \(a\) ; let \(\Lambda\) be the kernel of \(\exp_{A}\) , so that \(\Lambda\) is identified with \(a / \Lambda\) .
\end{enumerate}

\begin{enumerate}
\def\labelenumi{(\alph{enumi})}
\tightlist
\item
  The following conditions are equivalent:
\end{enumerate}

(al) The canonical mapping \(\mathbf{C} \otimes_{\mathbf{Z}} \Lambda \to a\) is injective.

(a2) A is isomorphic to a Lie subgroup of some \((\mathbf{C}^{*})^{n}\) .

(a3) A is isomorphic to a group \((\mathbf{C}^{*})^{p}\times \mathbf{C}^{q}\) .

(a4) A has a finite-dimensional faithful complex linear representation.

(a5) A has a finite-dimensional semi-simple faithful complex linear representation with closed image.

\begin{enumerate}
\def\labelenumi{(\alph{enumi})}
\setcounter{enumi}{1}
\tightlist
\item
  The following conditions are equivalent:
\end{enumerate}

(b1) The canonical mapping \(\mathbf{C} \otimes_{\mathbf{Z}} \Lambda \to \mathfrak{a}\) is surjective.

(b2) A is isomorphic to a quotient of a group \((\mathbf{C}^{*})^{n}\) .

(b3) No direct factor of A is isomorphic to C.

(b4) Every complex linear representation of A is semi-simple.

\begin{enumerate}
\def\labelenumi{(\alph{enumi})}
\setcounter{enumi}{2}
\tightlist
\item
  The following conditions are equivalent:
\end{enumerate}

\begin{enumerate}
\def\labelenumi{(\roman{enumi})}
\setcounter{enumi}{149}
\tightlist
\item
  The canonical mapping \(\mathbf{C} \otimes_{\mathbf{Z}} \Lambda \to \mathfrak{a}\) is bijective.
\end{enumerate}

(c2) A is isomorphic to a group \((\mathbf{C}^{*})^{n}\) .

\begin{enumerate}
\def\labelenumi{(\alph{enumi})}
\setcounter{enumi}{3}
\tightlist
\item
  Let \(\mathbf{F}\) be a finite subgroup of \(\mathbf{A}\) and let \(\mathbf{A}' = \mathbf{A} / \mathbf{F}\) . Show that \(\mathbf{A}\) satisfies conditions \((a_i)\) (resp. \((b_i), (c_i)\) ) if and only if \(\mathbf{A}'\) does.
\end{enumerate}

\begin{enumerate}
\def\labelenumi{\arabic{enumi}.}
\setcounter{enumi}{18}
\tightlist
\item
  Let \(G\) be a real Lie group. Show that there exists a neighbourhood \(V\) of 0 in \(L(G)\) such that, for \(x, y\) in \(V\) :
\end{enumerate}

\[
\exp (t (x + y)) = \lim _ {n \rightarrow + \infty} \left(\exp \frac {t x}{n}. \exp \frac {t y}{n}\right) ^ {n}
\]

\[
\exp (t ^ {2} [ x, y ]) = \lim _ {n \rightarrow + \infty} \left(\exp \frac {t x}{n}. \exp \frac {t y}{n}. \exp \frac {- t x}{n}. \exp \frac {- t y}{n}\right) ^ {n ^ {2}}
\]

uniformly for \(t \in (0, 1)\) .

\begin{enumerate}
\def\labelenumi{\arabic{enumi}.}
\setcounter{enumi}{19}
\item
  Let G be a Lie group, H a Lie subgroup of G and A an integral subgroup of G such that \(\mathrm{L}(\mathrm{H}) \cap \mathrm{L}(\mathrm{A}) = \{0\}\) . Show that \(H \cap A\) is discrete in the Lie group A.
\item
  Let \(G\) be a real Lie group and \(H\) a normal 1-dimensional integral subgroup.
\end{enumerate}

\begin{enumerate}
\def\labelenumi{(\alph{enumi})}
\item
  If H is not closed, \(\bar{H}\) is compact (Spectral Theories, Chapter II, §2, Lemma 1), hence isomorphic to \(T^{n}\) and hence central in G if G is connected (use General Topology, Chapter VII, §2, Proposition 5).
\item
  Suppose that H is closed. Let a be an element of G and C(a) the set of elements of G which commute with a. Suppose that H ⊕ C(a).
\end{enumerate}

If H is isomorphic to R, then \(\mathrm{C}(a) \cap \mathrm{H} = \{e\}\) . (Consider the automorphism \(\alpha: h \mapsto a^{-1}ha\) of H.) If H is isomorphic to T, \(\mathrm{C}(a) \cap \mathrm{H}\) has two elements. (Consider \(\alpha\) again and use General Topology, Chapter VII, §2, Proposition 6.) The second situation is impossible if G is connected (use \((a)\) ).

\begin{enumerate}
\def\labelenumi{(\alph{enumi})}
\setcounter{enumi}{2}
\tightlist
\item
  With the hypotheses of (b), suppose further that G/H is commutative.
\end{enumerate}

Then \(\mathrm{G} = \mathrm{C}(a)\) . H. (Observe that the mapping \(h \mapsto h^{-1}\alpha(h)\) of H into H is surjective.)

\begin{enumerate}
\def\labelenumi{\arabic{enumi}.}
\setcounter{enumi}{21}
\tightlist
\item
  Let G be a real Lie group, H a normal 1-dimensional integral subgroup and A a closed subgroup of G such that AH is not closed in G. Then H is central in the identity component of AH. (Reduce it to the case where G = AH and G is connected. Then the set B of elements of A which commute with H is a normal subgroup of G. Passing to the quotient by B, reduce it to the case B = \(\{e\}\) . As every commutator in G commutes with H, A is then commutative. By assuming that H is closed and using Exercise 21 (c), show that AH would be closed, whence a contradiction. Hence H is not closed and Exercise 21 (a) can be used.)
\end{enumerate}

¶ 23. Let G be a real Lie group and \(c = (\mathrm{U}, \phi, \mathrm{E})\) a chart on the manifold G centred on the identity element \(e\) . If \(x \in \mathrm{U}\) , let \(|x|_c\) denote the norm of \(\phi(x)\) in the Banach space E.

\begin{enumerate}
\def\labelenumi{(\alph{enumi})}
\tightlist
\item
  Show that, if \(c' = (\mathbf{U}', \phi', \mathrm{E}')\) is another chart centred on \(e\) , there exist constants \(\lambda > 0\) and \(\mu > 0\) such that
\end{enumerate}

\[
\mu | x | _ {c} \leqslant | x | _ {c ^ {\prime}} \leqslant \lambda | x | _ {c}
\]

for all \(x \in G\) sufficiently close to e.

\begin{enumerate}
\def\labelenumi{(\alph{enumi})}
\setcounter{enumi}{1}
\tightlist
\item
  Show that, for all \(\rho > 0\) , there exists a neighbourhood \(\mathbf{U}_{\rho}\) of \(e\) contained in \(\mathbf{U}\) such that, for \(x, y\) in \(\mathbf{U}_{\rho}\) , \((x, y) \in \mathbf{U}_{\rho}\) and
\end{enumerate}

\[
\left| (x, y) \right| _ {c} \leqslant \rho . \inf (\left| x \right| _ {c}, \left| y \right| _ {c}).
\]

\begin{enumerate}
\def\labelenumi{(\alph{enumi})}
\setcounter{enumi}{2}
\tightlist
\item
  Suppose that \(\mathbf{G}\) is finite-dimensional. Let \(\Gamma\) be a discrete subgroup of \(\mathbf{G}\) . Apply (b) with \(0 < \rho < 1\) and choose \(\mathrm{U}_{\rho}\) to be relatively compact. The set \(\mathrm{U}_{\rho} \cap \Gamma\) is finite. Let
\end{enumerate}

\[
\{e = \gamma_ {0}, \gamma_ {1}, \dots , \gamma_ {m} \}
\]

be its elements, numbered so that

\[
\left| \gamma_ {0} \right| _ {c} \leqslant \left| \gamma_ {1} \right| _ {c} \leqslant \dots \leqslant \left| \gamma_ {m} \right| _ {c}.
\]

Show that, if \(i \leqslant m\) and \(j \leqslant m\) , the commutator \((\gamma_{i}, \gamma_{j})\) is equal to one of the \(\gamma_{k}\) with \(k < \inf(i,j)\) . Deduce that the subgroup generated by \(\mathrm{U}_{\rho} \cap \Gamma\) is nilpotent of class \(\leqslant m\) .

Deduce from the above the existence of a neighbourhood \(\mathbf{V}\) of \(e\) such that, for every discrete subgroup \(\Gamma\) of \(\mathbf{G}\) , there exists a nilpotent integral subgroup \(\mathbf{N}\) of \(\mathbf{G}\) containing \(\mathbf{V} \cap \Gamma\) .

\begin{enumerate}
\def\labelenumi{(\alph{enumi})}
\setcounter{enumi}{3}
\tightlist
\item
  Suppose that G is connected and finite-dimensional and contains an increasing sequence of discrete subgroups \(D_{n}\) whose union is dense in G. Then G is nilpotent. (Using (c), prove the existence of a central one-parameter subgroup H which meets \(D_{n}\) for sufficiently large n at a point distinct from e. Argue by induction on dim G, distinguishing two cases according to whether H is closed or relatively compact (Exercise 21 (a)).
\end{enumerate}

\begin{enumerate}
\def\labelenumi{\arabic{enumi}.}
\setcounter{enumi}{23}
\tightlist
\item
  Let \(G, G'\) be finite-dimensional real Lie groups, \(f\) a surjective morphism of \(G\) into \(G'\) , \(N\) its kernel, \(H\) an integral subgroup of \(G\) and \(H' = f(H)\) .
\end{enumerate}

\begin{enumerate}
\def\labelenumi{(\alph{enumi})}
\item
  Suppose that N is finite. For \(H'\) to be closed, it is necessary and sufficient that H be closed.
\item
  Suppose that N is compact. If H is closed, \(H'\) is closed.
\end{enumerate}

\begin{enumerate}
\def\labelenumi{\arabic{enumi}.}
\setcounter{enumi}{24}
\tightlist
\item
  Let \(G\) be a finite-dimensional real or complex Lie group. Let \(S \subset L(G)\) . Suppose that \(L(G)\) is the Lie algebra generated by \(S\) and that \(S\) is stable under homotheties.
\end{enumerate}

\begin{enumerate}
\def\labelenumi{(\alph{enumi})}
\item
  Let H be the subgroup of G generated by exp S. Show that H is open in G. (Let A be the set of \(x \in L(G)\) such that \(\exp(Kx) \subset H\) and B the vector subspace of \(L(G)\) generated by A. Show that {[}A, S{]} \(\subset A\) and then that {[}A, A{]} \(\subset B\) . Deduce that B is a subalgebra of \(L(G)\) and then use Proposition 3 of § 4.)
\item
  Suppose further that
\end{enumerate}

\[
(x \in S \text {   and   } y \in S) \Rightarrow ((\operatorname{Ad} \exp x) (y) \in S).
\]

Then the vector subspace \(\mathbf{V}\) of \(\mathbf{L}(\mathbf{G})\) generated by \(\mathbf{S}\) is equal to \(\mathbf{L}(\mathbf{G})\) . (Using \((a)\) , show that \([\mathbf{L}(\mathbf{G}), \mathbf{V}] \subset \mathbf{V}\) .)

\begin{enumerate}
\def\labelenumi{\arabic{enumi}.}
\setcounter{enumi}{25}
\tightlist
\item
  Let G be an n-dimensional connected compact complex Lie group, X a finite-dimensional connected complex analytic manifold and \((g, x) \mapsto gx\) a law of analytic left operation of G on X. For all \(g \in G\) , let \(\rho(g)\) be the mapping \(x \mapsto gx\) of X into X. Suppose, that, for all \(g \in G\) distinct from \(e, \rho(g) \neq Id_x\) . Then every orbit of G in X is an n-dimensional closed submanifold of X. (Let \(x \in X\) , H be the stabilizer of x and \(H'\) the identity component of H. For all \(g \in H'\) , let \(u(g)\) be the tangent mapping at x to \(\rho(g)\) . Arguing as in Proposition 6 of no. 3, show that \(u(g)\) is the identity for all \(g \in H'\) . Using § 1, Exercise 4 and the connectedness of X, deduce that \(H' = \{e\}\) .)
\end{enumerate}

¶ 27. Let G be a compact real Lie group, M a compact manifold of class C² and J an open interval of R containing 0. Let \((s, \langle x, \xi \rangle) \mapsto (m_{\xi}(s, x), \xi)\) be a mapping of class C² of G × (M × J) into M × J under which G operates on M × J on the left.

\begin{enumerate}
\def\labelenumi{(\alph{enumi})}
\item
  Let X be a vector field of class \(C^{1}\) on \(M \times J\) such that, for all \((x, \xi) \in M \times J\) , the second projection of \(\mathbf{X}_{(x, \xi)}\) is the tangent vector l to J. Translating X by G and integrating over G, deduce the existence of a field \(X'\) with the same properties as X, which is moreover invariant under G.
\item
  Show that there exists a diffeomorphism \((x, \xi) \mapsto (h_{\xi}(x), \xi)\) of \(\mathbf{M} \times \mathbf{J}\) onto itself such that
\item
  for all \(\xi \in \mathbf{J}\) , \(h_{\xi}\) is a diffeomorphism of \(\mathbf{M}\) onto \(\mathbf{M}\) ;
\end{enumerate}

\begin{enumerate}
\def\labelenumi{(\roman{enumi})}
\setcounter{enumi}{1}
\tightlist
\item
  \(m_{\xi}(s,x) = h_{\xi}(m_0(s,h_{\xi}^{-1}(x)))\) for \(s\in \mathbf{G}, x\in \mathbf{M},\xi \in \mathbf{J}\) .
\end{enumerate}

(Use (a) and Theorem 5 of no. 8.)

\begin{enumerate}
\def\labelenumi{\arabic{enumi}.}
\setcounter{enumi}{27}
\item
  Give an example of a finite-dimensional real Lie group G and two integral subgroups A, B of G such that \(\mathrm{L}(\mathrm{G})=\mathrm{L}(\mathrm{A})+\mathrm{L}(\mathrm{B})\) but AB is not a subgroup of G (cf.~Integration, Chapter VII, § 3, no. 3, Proposition 6).
\item
  Let G be a finite-dimensional complex connected Lie group, \(G_{0}\) the underlying real Lie group and H an integral subgroup of \(G_{0}\) . Show that there exists a smallest integral subgroup \(H^{*}\) of G containing H. Give an example where H is closed in \(G_{0}\) but \(H^{*}\) is not closed in G (take \(G = C^{2}/Z^{2}\) ).
\item
  Let \(G\) be a compact connected complex Lie group and hence of the form \(\mathbf{C}^n / D\) , where \(D\) is a discrete subgroup of \(\mathbf{C}^n\) of rank \(2n\) .
\end{enumerate}

\begin{enumerate}
\def\labelenumi{(\alph{enumi})}
\item
  Show that every holomorphic differential 1-form \(\omega\) on G is invariant. (Let \(\pi: \mathbf{C}^n \to \mathbf{G}\) be the canonical morphism. Let \(\zeta_1, \ldots, \zeta_n\) be the coordinate functions on \(\mathbf{C}^n\) . Then \(\pi^*(\omega)\) is of the form \(\sum_j a_j d\zeta_{j}\) , where the \(a_j\) are holomorphic functions on \(\mathbf{C}^n\) which are invariant under D and hence constant.)
\item
  Let \(G' = C^{n}/D'\) , where \(D'\) is a discrete subgroup of \(C^{n}\) of rank 2n. Show that every analytic manifold isomorphism \(u: G \to G'\) is of the form \(s \mapsto v(s) + a'\) , where v is a Lie group isomorphism and \(a' \in G'\) . (Let \(\pi': C^{n} \to G'\) be the canonical morphism. There exists an isomorphism of analytic manifolds \(\tilde{u}: C^{n} \to C^{n}\) such that \(u \circ \pi' \circ \tilde{u}\) . For every holomorphic differential 1-form \(\omega'\) on \(G'\) , \(u^*(\pi^*(\omega'))\) is a 1-form on \(C^{n}\) which is invariant under translations by (a). Deduce that \(\tilde{u}\) is an affine mapping.)
\end{enumerate}

